% PNP RESEARCH REPORT::RECOVERY-STRENGTHEN6::2026-10-08T22:53:00Z
\documentclass[11pt]{article}
\usepackage[margin=0.85in]{geometry}
\usepackage{amsmath,amssymb,amsthm,booktabs,hyperref,xurl}
\emergencystretch=3em
\newtheorem{theorem}{Theorem}
\newtheorem{proposition}[theorem]{Proposition}
\newcommand{\SAT}{\operatorname{SAT}}
\newcommand{\MCB}{\mathsf{MCB}}
\title{RECOVERY-STRENGTHEN6:\\What the Earlier WRAPPERs Actually Close\\Hereditary-Certificate $K$ Elimination and the Flat-Child Obstruction}
\author{Recovery proof-strengthening agent}
\date{8 October 2026 UTC}

\begin{document}\sloppy\maketitle

\paragraph{Status.}
\textsc{Proved/Exact}: a corrected recovery classification, a promise-class
successor-closure theorem, and an asymptotic counterfamily showing that the
current flat marked compiler is not recursively closed.  \textsc{Computational}:
exact finite CNF1 checks of the counterfamily and exact rational checks of the
recovered recurrence and current FLMV margin.  Earlier work did build one
terminating exact source-backed solver on every finite CNF; earlier work also
proved conditional recursive elimination of the inherited exponent $K$.
Neither fact proves that one implementation has the favorable all-input
time, child-volume and simultaneous-space bounds.  P versus NP remains
\textsc{Open}.  The retired unrestricted $t+s<2\cos(\pi/7)$ endpoint is not
used.

\section{Claim card and exact quantifiers}
Fix one encoded SAT language over all bit strings, with encoded bit length
$N$, and one deterministic random-access machine realization.  The
load-bearing target under the counterfactual $\SAT\in P$ is
\[
 \exists R_A\;\exists h<1\;\exists r,\eta,p\quad
 \forall x\in\mathrm{SATEnc}\quad
 \forall y\text{ recursively reached from }x,                 \tag{1}
\]
such that the same uniform source-first procedure either decides the current
instance within the displayed local ledger or emits exact ordinary-SAT exits,
with all discovery, failed attempts, multiplicity, output, replay, buffers,
indexes and ancestors charged.  The favorable concrete contract currently
uses
\[
 h=\frac12,
 \qquad \eta=\frac35,
 \qquad r\le\frac65,
 \qquad p\le\frac12.                                         \tag{2}
\]
One size-only $o(1)$ envelope must hold uniformly across inputs, true adaptive
transcripts and recursion depth.  Time and space must be simultaneous bounds
for the same machine.

MARKED-CUTSET-BLOCK1 instead proves
\[
 \exists R\quad\forall x\in D_{\rm marked},                  \tag{3}
\]
where $D_{\rm marked}$ is a supplied annotated subclass.  Its emitted children
are ordinary CNFs, not guaranteed members of $D_{\rm marked}$.  The difference
between (1) and (3) is the remaining recursive-closure/acquisition obligation.

\section{Recovery correction: what was already implemented}
The older corpus contains more implementation than the recent shorthand
``source-backed replacement is still unbuilt'' allowed.

\begin{center}
\begin{tabular}{p{.25\linewidth}p{.31\linewidth}p{.35\linewidth}}
\toprule
artifact & exact achievement & missing favorable premise\\
\midrule
COMBINED-REPLAY1 & one terminating exact solver for every finite CNF using source-backed states, frontier DP, unit steps, component AND, Shannon OR and collision-checked bounded sharing & linear clause scopes/maps may remain, source replay is parent-dependent, and Shannon multiplicity may be exponential\\
TRIE-BATCHACQUIRE1 & exact shared residual-delta object from one source scan for a supplied finite assignment trie & parent-aware object is not a stand-alone small ordinary CNF; output, recovery and fanout remain charged\\
RECURCLOSE1 & $C=\max\{r,\eta/(1-\theta)\}$ and depth-first space closure, independent of the initial exponent $K$ & assumes one uniform all-input strict shrinker\\
CYCLEGAIN1 & exact typed closure with well-founded same-scale internals and contraction on every complete return cycle & assumes the source-matched typed graph and low-space interfaces\\
STREAMCOMP1 & exact per-access time/space charge for a supplied random-access child generator & does not construct the generator\\
ORBITBATCH1 & exact field-address deduplication and shared setup in one SU05/Hadamard realization & does not deduplicate executions or outputs and is not normalized to SAT encoded length\\
\bottomrule
\end{tabular}
\end{center}

Thus three statements are independently true:
\[
 \begin{array}{ll}
 \text{exact all-input source-backed implementation} & \textbf{YES},\\
 \text{conditional recursive $K$ elimination} & \textbf{YES},\\
 \text{one implementation satisfying every premise in (1)--(2)} & \textbf{NOT PROVED}.
 \end{array}                                               \tag{4}
\]
This corrects a stale implementation summary without promoting any resource
claim.

\section{Recovered promise-class closure}
Call a domain $D$ \emph{hereditary for $\MCB$} if a single source-checkable
encoding grammar has the following property at every sufficiently large
encoded length $N$: every $x\in D$ is accepted by the source-first marked
phase, and every complete encoded child emitted on every true transcript is
again in $D$.  Any nested certificate bits, local renumbering map, child
header and verification work are included in the child's encoded length and
the parent's local ledger.  No semantic SAT call may be used to establish
membership.  The same grammar, verifier and constants apply at every level.

\begin{theorem}[Hereditary marked-certificate closure]
Suppose $D$ is hereditary and one uniform depth-first realization, on every
$x\in D$ of length $N$, has
\[
 T_{\rm local}(N)\le N^{11/10+o(1)},\quad
 q(N)\le N^{3/5+o(1)},\quad
 \max_i |y_i|\le N^{1/2+o(1)},                              \tag{5}
\]
uses $N^{1/2+o(1)}$ peak local state including its one live child, and is
exact under recursive child answers.  Then the same recursive machine decides
$D$ in
\[
 T(N)=N^{6/5+o(1)},\qquad S(N)=N^{1/2+o(1)}.                \tag{6}
\]
The bound is independent of the exponent of any original polynomial SAT
decider.
\end{theorem}

\begin{proof}
Hereditary closure makes the same verified phase applicable to every
descendant; strict power shrink reaches a fixed finite base.  RECURCLOSE1
applied to (5) gives
\[
 C=\max\left\{\frac{11}{10},
   \frac{3/5}{1-1/2}\right\}=\frac65.                       \tag{7}
\]
Repeated equal children remain repeated executions in $q$.  A depth-first
machine retains one fully charged frame and one child at each level.  Summing
the power-shrinking frames gives exponent
$\max\{1/2,1/2\}=1/2$.  This is one joint realization, not a splice of a
fast machine and a low-space machine.  Uniform error envelopes give the
standard strong-induction proof after arbitrary fixed exponent slack.
\end{proof}

The theorem identifies exactly what a genuinely recursive annotation would
buy.  It is a promise-class theorem.  No current artifact constructs such a
hereditary certificate for every general CNF or even implements the proposed
nested tag grammar and local renumbering.  Those steps remain candidates, not
proved facts.

\section{PROVED/EXACT: the current flat language is not hereditary}
Let $G_m$ be any growing family of well-formed unmarked CNFs, using variables
at least $4$, with exact CNF1 encoded length $m$.  Construct a marked parent
$P_m$ containing the magic marker, cutset $S=\{3\}$, the cut terminator,
$R_m$ copies of the global tautology $(3\vee\neg3)$, one interval marker
covering $G_m$, and then the clauses of $G_m$.  Take
$R_m=\Theta(m^2)$.

\begin{proposition}[Flat-child obstruction]
For all sufficiently large $m$, $P_m$ is accepted by the current flat marked
compiler, and both cut assignments emit the identical child $G_m$.  The child
is outside the marked domain, so the next invocation returns UNKNOWN.  Hence
the current flat domain fails successor closure beyond every fixed finite base.
\end{proposition}

\begin{proof}
All added global clauses are tautologies over the cutset, so they disappear
under either assignment and do not change satisfiability.  The complete parent
length is $N=\Theta(m^2)$; all padding is charged.  For sufficiently large
$N$, $k=1\le\lfloor\log_2N/10\rfloor$, there is one block, identifiers obey
the $N^2$ cap, and
\[
 |G_m|=m\le \sqrt N\log_2N.                                \tag{8}
\]
Thus the dry run accepts and executes two child occurrences.  The emitted
stream begins with $G_m$'s ordinary CNF header and lacks the magic marker, so
the same parser returns UNKNOWN.  Since $m$ is unbounded, no fixed base table
repairs this failure.  \qedhere
\end{proof}

This does not challenge one-level correctness.  It falsifies only the
promotion from the flat partial compiler to the all-level premise of
RECURCLOSE1.

\section{Why a trivial double wrapper does not close the cycle}
Use two types: $U$ for an ordinary unmarked $m$-bit CNF and $M$ for a marked
instance.  The current compiler provides $M\to U$ with scale exponent
$1/2+o(1)$ only after a valid annotation has been supplied.  A contracting
typed cycle additionally needs a uniform $U\to M$ transition.

Putting an arbitrary $U$ instance into one marked block with only linear
overhead fails (8): its only residual child still has $\Theta(m)$ bits, larger
than $\sqrt{\Theta(m)}\log m$.  Padding until the flat parser accepts requires
\[
 N=\Omega\!\left(\frac{m^2}{\log^2m}\right).                \tag{9}
\]
That $U\to M$ edge expands with exponent $2-o(1)$, outside CYCLEGAIN1's
nonexpanding edge model.  The complete macro returns an ordinary child of
size $m$; its exponent product is $(2-o(1))(1/2+o(1))=1+o(1)$, not a proved
strict contraction.  Therefore a padding wrapper plus the marked wrapper is
not the needed double wrapper.  A genuinely decomposing annotation is exactly
the missing acquisition construction.

Falling back to the assumed decider $A$ on $U$ is exact but restores the term
$\tfrac12K$ in time and potentially uncontrolled workspace.  CYCLEGAIN1
cannot manufacture the missing $U\to M$ edge.  STREAMCOMP1 charges a supplied
implicit edge but does not create it.  ORBITBATCH1 deduplicates field addresses,
not the two executed copies of $G_m$; avoiding the second execution would need
a separately implemented and charged exact cache rule.

\section{Primary-literature check}
M\"ahlmann, Siebertz and Vigny formalize recursive backdoors with assignment
(OR) and component (AND) nodes.  Their supplied-tree semantics and their
$2^k|F|$ leaf/occurrence bound validate the logical structure of a hereditary
certificate.  Their detection theorem is parameterized: in time
$2^{2^{O(k)}}|G|$ it returns a recursive backdoor of bounded parameterized
depth or proves that the optimum depth exceeds $k$.  This is useful on a
bounded-parameter promise; for $k=\Theta(\log N)$ it does not provide the
specific $N^{6/5}$ all-input acquisition ledger:
\url{https://doi.org/10.4230/LIPIcs.MFCS.2021.73}.

Dreier, Ordyniak and Szeider extend bounded-depth backdoors for Horn and Krom
and obtain linear-time SAT on each fixed bounded-depth class.  Their FPT
approximation dependence is not an all-input power-shrink theorem and they do
not state the required same-machine half-space bound:
\url{https://doi.org/10.4230/LIPIcs.ESA.2022.46}.

Known treewidth acquisition results display the joint-resource issue rather
than resolving it.  Kammer, Meintrup and Sajenko's Theorem~31 and
Corollary~32 output approximate decompositions in $c^kn$ times polylogarithmic
factors while using $O(n+k^2\log^2n)$ bits.  That is linear, not half-space,
in the sparse regime.  Their introduction notes that a cited $o(n)$-bit
alternative costs $\Omega(n^2)$ time:
\url{https://arxiv.org/abs/1907.00676}.  One may not combine the time theorem
from one machine with the space theorem from another.

These papers support
\[
 \text{supplied recursive certificate}\Longrightarrow
 \text{exact recursive replay},                             \tag{10}
\]
not
\[
 \SAT\in P\Longrightarrow
 \text{cheap hereditary-certificate acquisition on every input}. \tag{11}
\]

\section{Endpoint and encoding guard}
The current finite-entry check uses the 2005 Fortnow--Lipton--van
Melkebeek--Viglas deterministic random-access model, with work and index tapes
charged.  Corollary~5.8 gives
\[
 d(c)=\frac12\left(\frac{c+2}{c^2}-c\right).
\]
At $c=1201/1000$, exact arithmetic gives
\[
 d(c)=\frac{1468676399}{2884802000}>\frac12,
 \qquad c>\frac65.                                         \tag{12}
\]
Thus a correctly transferred same-machine $(6/5,1/2)$ SAT algorithm would
enter the strict primary forbidden region.  The source is
\url{https://lance.fortnow.com/papers/files/tst.pdf}, especially its model in
Section~2.1, local transfer in Theorem~3.1, and Corollary~5.8.

The local encoding transfer is still open for the project's exact CNF1
syntax.  Tautological annotations may be deleted by ordinary normalizers.
Any endpoint use must define malformed inputs and locally generate the target
encoding with polylogarithmic time and logarithmic workspace per requested
bit.  Materializing a whole near-linear reduction would violate half-space.

\section{Bounded reproducible validation}
The companion program
\path{experiments/proof_spines/hereditary_marker_closure_test.py} has SHA-256
\nolinkurl{e893ff9b0c584525d8f3c42a0fb1e2b692f04c2131babd16edd2aea1b5527934}.
With a 25-second alarm it checked six increasing exact CNF1 families.  Every
padded parent was accepted, executed the same ordinary child twice, obeyed
the exact square-root/log child cap, and produced a successor rejected for a
missing magic marker.  The largest case had 401,624 parent bits and an 8,248
bit child.  It also checked (7) and (12) with exact rational arithmetic.  The
run passed in 0.13 seconds.  The result file is
\path{experiments/proof_spines/hereditary_marker_closure_test.json}.  These are
finite parser/arithmetic checks; Python memory is not bit-machine space
evidence and the test does not implement a hereditary grammar.

\section{Corpus, reconciliation and handoff}
\paragraph{Corpus reviewed.}
Current main4/main5, root status, AGENTS, CODEX, the repository research skill
and protocol, FORMAL-SPINE1, latest Construction/Recovery/ARC/Red-Team reports,
the exact original RECURCLOSE1, CYCLEGAIN1, STREAMCOMP1 and ORBITBATCH1 blobs,
COMBINED-REPLAY1, REPLAY-RECOVERY, TRIE-BATCHACQUIRE1 and later replay/index
records, plus the cited primary literature.  Original blob identifiers are
respectively
\nolinkurl{5c77ed64e79493883a52d0b5425dd195ba17b3cf},
\nolinkurl{aa3859c364c0c6eb15a653f0218105cf9eb1917a},
\nolinkurl{adc67bbc9f5d0255defc770cd48117dece0741d7}, and
\nolinkurl{766891e34f5151ed542e407d0f8bfee63322c6ba}.

\paragraph{Reconciliation.}
ADOPT the correction that exact source-backed all-input replay was built.
ADOPT the hereditary promise-class closure theorem as the exact way the marked
phase and RECURCLOSE1 would compose.  REJECT recursive promotion of the
current flat compiler.  PRESERVE multiplicity, replay and simultaneous-space
charges, conditional $K$ elimination, the FLMV strict endpoint and explicit
OPEN status.

\paragraph{Next target.}
Either implement and verify a nested successor-closed marker grammar and then
separately acquire it causally on every ordinary CNF, or strengthen
COMBINED-REPLAY1 itself by proving an all-input completed-phase bound.  For the
latter, the exact missing estimates are aggregate executed child work/volume,
fixed-power largest descendant, and simultaneous peak state including source
replay scopes and ancestors.  The first regression test for any marked route
is the flat-child family above; the first general-input hostile control remains
the expander-guarded core from CONSTRUCTION-STRENGTHEN6.

\paragraph{Final status.}
The earlier wrappers solved exact replay and conditional accounting, and the
new marked wrapper supplies a favorable one-level causal member.  Their valid
composition currently stops at a supplied hereditary promise.  No reviewed
artifact proves the all-input favorable implementation or the implication
(11).  P versus NP remains \textsc{Open}.

% END_OF_VERIFIED_REPORT::RECOVERY-STRENGTHEN::RECOVERY-STRENGTHEN6
\end{document}

